\chapter*{TÓM TẮT}
\addcontentsline{toc}{chapter}{TÓM TẮT}

Đồ án đánh giá VPC Peering và AWS Transit Gateway (TGW) trên cùng bốn
endpoint EC2 nhằm tách ảnh hưởng của đường kết nối khỏi khác biệt cấu hình
máy. Hạ tầng gồm bốn VPC trong \texttt{us-east-1}, mỗi VPC có một EC2
\texttt{c6i.large} tại cùng Availability Zone. Terraform giữ nguyên endpoint
và chuyển 12 route liên VPC giữa full-mesh Peering và TGW hub-and-spoke.
Ba bài đo chính gồm RTT khi mạng rảnh, TCP một luồng từ A tới B và TCP fan-in
từ A, C, D tới B. Mỗi bộ dữ liệu có năm lượt; raw data, metadata, CPU và bộ
đếm ENA được lưu để kiểm chứng lại bằng script.

Trong các phiên đo của nghiên cứu này, RTT trung bình là 0,2078 ms qua
Peering và 0,5630 ms qua TGW. Receiver goodput TCP một luồng đạt lần lượt
4,7828 và 4,6083 Gbps. Tổng fan-in xấp xỉ đạt 11,6110 và 11,5905 Gbps,
nhưng cả hai phiên đều ghi nhận bộ đếm network allowance tại máy nhận tăng.
Kết quả fan-in vì vậy phản ánh đồng thời đường mạng và giới hạn endpoint,
không chứng minh hai dịch vụ có năng lực tối đa tương đương. Báo cáo kết luận
Peering có lợi thế về latency và single-flow trong cấu hình đã đo, còn TGW có
lợi thế về topology, định tuyến tập trung và khả năng mở rộng quản trị. Phạm
vi kết luận bị giới hạn bởi số lượt nhỏ, phép đo hai mode nối tiếp, đặc tính
burst của EC2 và cửa sổ fan-in chưa đồng bộ tuyệt đối.
